\chapter{\texttt{with-eval-after-load} vs \texttt{:after}}
\label{ch:with-eval-after-load}
\minitoc\label{with-eval-after-load-vs-after}
\index{with-eval-after-load}\index{after-load-alist}\index{feature!provide}\index{load ordering}

% ── Opening signpost ──────────────────────────────────────────────
Chapter~\ref{ch:use-package} introduced the \texttt{:after} keyword;
Emacs also provides a lower-level mechanism,
\texttt{with-eval-after-load}, for running code after a package loads.
This chapter compares the two approaches and shows when to prefer
each one.

\texttt{(info "(elisp) Hooks for Loading")}

The config calls \texttt{(with-eval-after-load 'eglot …)}, \texttt{(with-eval-after-load 'which-key …)}, and \texttt{(with-eval-after-load 'cc-mode …)} at top-level outside any \texttt{use-package} block. These look similar to \texttt{:after} but have fundamentally different semantics.

\section{How \texttt{with-eval-after-load} works}\label{how-with-eval-after-load-works}

\texttt{with-eval-after-load} queues a form in the internal alist \texttt{after-load-alist} under a feature key. When any file is loaded and calls \texttt{(provide 'FEATURE)}, Emacs drains the queue for that feature and evaluates every accumulated form.

\begin{minted}{elisp}
;; Queues the lambda under the key 'eglot:
(with-eval-after-load 'eglot
  (add-to-list 'eglot-server-programs
               '(web-mode . ("typescript-language-server" "--stdio"))))
\end{minted}

If the feature has \textbf{already been loaded} when \texttt{with-eval-after-load} is called, the form is evaluated \textbf{immediately}. If not, it waits.

The internal structure of \texttt{after-load-alist} (simplified):

\begin{minted}{elisp}
;; ((eglot  . ((lambda () (add-to-list 'eglot-server-programs …))))
;;  (which-key . ((lambda () (which-key-add-key-based-replacements …)))))
\end{minted}

\section{How \texttt{:after} in \texttt{use-package} works}\label{how-after-in-use-package-works}

\texttt{:after pkg} is a \textbf{load-order gate}: it prevents the enclosing \texttt{use-package} body from being processed until \texttt{pkg} has been loaded. Under the hood it wraps the body in \texttt{with-eval-after-load}. Combined with \texttt{:demand t} it means "load this package immediately once its dependency loads".

\section{When to use each}\label{when-to-use-each}

Table~\ref{tab:eval-after-load-situations} compares common situations where
you would reach for one mechanism or the other.

\begin{sidewaystable}[p]
\centering
\caption{When to use \texttt{with-eval-after-load} vs.\ \texttt{:after}}
\label{tab:eval-after-load-situations}
\begin{tabularx}{\linewidth}{@{}Xl@{}}
\toprule
Situation & Mechanism \\
\midrule
Registering a server in \texttt{eglot-server-programs} from a top-level form & \texttt{with-eval-after-load 'eglot} \\
Defining keybindings in a mode-map provided by package A, declared inside \texttt{use-package} for B & \texttt{:after A} in B's declaration \\
Configuration that must run after \textbf{two} packages both load & \texttt{:after (A B)} \\
Keymap setup for a built-in package outside any \texttt{use-package} & \texttt{with-eval-after-load 'package} \\
\bottomrule
\end{tabularx}
\end{sidewaystable}

\section{\texttt{with-eval-after-load} does not trigger loading}\label{with-eval-after-load-does-not-trigger-loading}

Unlike \texttt{:hook}, \texttt{:bind}, or \texttt{:commands}, \texttt{with-eval-after-load} does \textbf{not} cause a package to load. If nothing else triggers \texttt{eglot} to load, the queued form never runs. The pattern works here because \texttt{eglot-ensure} is on \texttt{prog-mode-hook}, guaranteeing \texttt{eglot} is always loaded before using its server-programs alist.

% ── Mental model summary ────────────────────────────────────────────
\section{Mental model}
\texttt{with-eval-after-load} queues a form in the internal \texttt{after-load-alist}; when the named feature's \texttt{provide} eventually executes, the queued form runs.  If the feature is already loaded at the time the \texttt{with-eval-after-load} is evaluated, the form runs immediately instead of being deferred.  Inside \texttt{use-package}, the \texttt{:after} keyword provides the same mechanism: it wraps the entire body in \texttt{with-eval-after-load}, acting as a load-order gate that ensures the body only runs once its dependencies are present.

Crucially, neither \texttt{with-eval-after-load} nor \texttt{:after} \emph{triggers} loading; something else, such as a hook, a keybinding, or an explicit \texttt{require}, must cause the feature to load before the queued form will fire.  In practice, use \texttt{with-eval-after-load} for top-level configuration outside \texttt{use-package} blocks, and use \texttt{:after} inside \texttt{use-package} for expressing load-order dependencies between packages.
